% =============================================================================
% paper46_galoistheorie_de.tex
% Paper 46: Galois-Theorie — Von Körpererweiterungen zur Unlösbarkeit der Quintik
% Specialist Mathematics Project — Build 135
% Autor: Michael Fuhrmann
% Datum: 2026-03-12
% =============================================================================

\documentclass[12pt,a4paper]{amsart}
\usepackage{amsmath,amssymb,amsthm}
\usepackage{mathtools}
\usepackage[utf8]{inputenc}
\usepackage[T1]{fontenc}
\usepackage[ngerman]{babel}
\usepackage{hyperref}
\usepackage{enumitem,booktabs,array}
\usepackage{geometry}
\geometry{margin=2.5cm}

% ---------------------------------------------------------------------------
% Theorem-Umgebungen
% ---------------------------------------------------------------------------
\newtheorem{theorem}{Satz}[section]
\newtheorem{lemma}[theorem]{Lemma}
\newtheorem{korollar}[theorem]{Korollar}
\newtheorem{proposition}[theorem]{Proposition}
\newtheorem{vermutung}[theorem]{Vermutung}
\theoremstyle{definition}
\newtheorem{definition}[theorem]{Definition}
\newtheorem{beispiel}[theorem]{Beispiel}
\newtheorem{bemerkung}[theorem]{Bemerkung}

% ---------------------------------------------------------------------------
% Metadaten
% ---------------------------------------------------------------------------
\title{Galois-Theorie:\\
       Von Körpererweiterungen zur Unlösbarkeit der Quintik}

\author{Michael Fuhrmann}
\address{Specialist Mathematics Project, Build 135}
\email{specialist-maths@localhost}
\date{\today}

\keywords{Galois-Theorie, Körpererweiterungen, Zerfällungskörper, Galois-Gruppe,
          Hauptsatz, Lösbarkeit durch Radikale, Satz von Abel-Ruffini,
          Kronecker-Weber-Satz, Konstruierbarkeit}

\subjclass[2020]{12F10, 12F20, 11R32, 12E05}

% ---------------------------------------------------------------------------
% Makros
% ---------------------------------------------------------------------------
\newcommand{\Q}{\mathbb{Q}}
\newcommand{\R}{\mathbb{R}}
\newcommand{\C}{\mathbb{C}}
\newcommand{\Z}{\mathbb{Z}}
\newcommand{\F}{\mathbb{F}}
\newcommand{\Gal}{\operatorname{Gal}}
\newcommand{\Aut}{\operatorname{Aut}}
\newcommand{\ord}{\operatorname{ord}}
\newcommand{\Fix}{\operatorname{Fix}}
\newcommand{\chara}{\operatorname{char}}

% =============================================================================
\begin{document}
% =============================================================================

\begin{abstract}
Die Galois-Theorie gehört zu den tiefgründigsten Errungenschaften der
Mathematikgeschichte. Von \'{E}variste Galois 1832 entwickelt und posthum
veröffentlicht, stellt sie eine tiefe Korrespondenz zwischen den Symmetriegruppen
der Polynomwurzeln und der Arithmetik von Körpererweiterungen her. Diese Arbeit
präsentiert eine in sich geschlossene Einführung in die Galois-Theorie: Ausgehend
von den Grundlagen der Körpererweiterungen und algebraischen Abschlüsse über
Zerfällungskörper, Separabilität, Normalität und Galois-Gruppe bis hin zum
berühmten Hauptsatz der Galois-Theorie. Anschließend wird diese Maschinerie
angewandt, um die klassischen Probleme der Lösbarkeit durch Radikale — insbesondere
der Satz von Abel-Ruffini — und die Konstruierbarkeit mit Zirkel und Lineal zu
klären. Der Kronecker-Weber-Satz über abelsche Erweiterungen von $\Q$ wird als
weiteres bedeutendes Anwendungsbeispiel vorgestellt. Durchgehend illustrieren
konkrete Beispiele die abstrakte Theorie.
\end{abstract}

\maketitle
\tableofcontents

% =============================================================================
\section{Einleitung}
% =============================================================================

Das Problem, Polynomgleichungen durch Radikale zu lösen — also durch
\emph{geschlossene Formeln} mit den Koeffizienten und verschachtelten Wurzeln —
beschäftigte Mathematiker über Jahrhunderte. Die Babylonier kannten die
Lösungsformel für quadratische Gleichungen um 2000 v.\,Chr.; Cardano
veröffentlichte die kubische Formel in der \emph{Ars Magna} (1545), und
Ferrari leitete im selben Jahrzehnt eine Formel für Grad vier her. Aber
Grad fünf widerstand allen Anstrengungen bis 1799, als Paolo Ruffini
(mit einem unvollständigen Beweis) ankündigte, keine solche Formel könne
existieren; Abel lieferte 1824 den ersten rigorosen Beweis.

Die entgültige Erklärung kam von \'{E}variste Galois (1811--1832), der in
einem Memoire, das er in der Nacht vor seinem tödlichen Duell schrieb, erschuf,
was wir heute \emph{Galois-Theorie} nennen. Galois ordnete jedem Polynom
$f \in K[x]$ eine Symmetriegruppe $\Gal(f/K)$ seiner Wurzeln zu und bewies:
\[
  f \text{ ist durch Radikale lösbar} \iff \Gal(f/K) \text{ ist eine auflösbare Gruppe.}
\]
Da das generische Grad-5-Polynom die Galois-Gruppe $S_5$ besitzt und
$A_5 \trianglelefteq S_5$ einfach und nicht-abelsch ist, ist $S_5$ nicht
auflösbar; daher existiert keine Radikalformel für Grad fünf.

Die Theorie hat sich weit über ihren Ursprung hinaus entwickelt. Heute ist
die Galois-Theorie zentral für die algebraische Zahlentheorie, das
Langlands-Programm und die moderne Kryptographie.

\subsection*{Historische Zeittafel}

\begin{center}
\begin{tabular}{ll}
\toprule
Jahr & Entwicklung \\
\midrule
$\sim$2000 v.\,Chr. & Babylonische Lösungsformel für quadratische Gleichungen \\
1545 & Cardano/Ferrari: kubische und quartische Formeln (\emph{Ars Magna}) \\
1799 & Ruffini: erster (unvollständiger) Beweis der Unlösbarkeit für Grad 5 \\
1824 & Abel: strenger Beweis des Satzes von Abel-Ruffini \\
1832 & Galois: gruppentheoretisches Kriterium; Memoire posthum veröffentlicht \\
1870 & Jordan: erstes Buch über Gruppentheorie mit Galois-Theorie \\
1801/1853 & Gauß/Kronecker-Weber: Kreisteilungskörper und abelsche Erweiterungen \\
\bottomrule
\end{tabular}
\end{center}

% =============================================================================
\section{Körpererweiterungen}
\label{sec:koerper}
% =============================================================================

\subsection{Grundlegende Definitionen}

\begin{definition}
Eine \emph{Körpererweiterung} $L/K$ ist eine Einbettung von Körpern
$K \subseteq L$. Der \emph{Grad} $[L:K] = \dim_K L$ ist die Dimension von
$L$ als $K$-Vektorraum. Die Erweiterung heißt \emph{endlich}, wenn
$[L:K] < \infty$.
\end{definition}

\begin{theorem}[Turmsatz]
Sind $K \subseteq F \subseteq L$ Körper, so gilt
\[
  [L:K] = [L:F] \cdot [F:K].
\]
\end{theorem}

\begin{proof}
Sei $\{e_i\}$ eine Basis von $F$ über $K$ und $\{f_j\}$ eine Basis von $L$
über $F$. Man zeigt, dass $\{e_i f_j\}$ eine $K$-Basis von $L$ bildet, also
$[L:K] = |I| \cdot |J| = [F:K] \cdot [L:F]$.
\end{proof}

\subsection{Algebraische und transzendente Elemente}

\begin{definition}
Ein Element $\alpha \in L$ heißt \emph{algebraisch über $K$}, wenn ein
nichttriviales Polynom $f \in K[x]$ mit $f(\alpha) = 0$ existiert. Andernfalls
heißt $\alpha$ \emph{transzendent über $K$}. Die Erweiterung $L/K$ heißt
\emph{algebraisch}, wenn jedes Element von $L$ algebraisch über $K$ ist.
\end{definition}

\begin{definition}
Das \emph{Minimalpolynom} $\mathrm{min}_{K}(\alpha)$ eines algebraischen
Elements $\alpha$ über $K$ ist das eindeutige normierte irreduzible Polynom
in $K[x]$, das $\alpha$ als Nullstelle hat.
\end{definition}

\begin{proposition}
Ist $\alpha$ algebraisch über $K$ mit Minimalpolynom $p$ vom Grad $n$, so gilt
$K(\alpha) \cong K[x]/(p)$ und $[K(\alpha):K] = n$.
\end{proposition}

\begin{beispiel}
Das Minimalpolynom von $\sqrt{2}$ über $\Q$ ist $x^2 - 2$, also
$[\Q(\sqrt{2}):\Q] = 2$ und $\Q(\sqrt{2}) = \{a + b\sqrt{2} : a,b \in \Q\}$.
Entsprechend gilt $\mathrm{min}_\Q(\sqrt[3]{2}) = x^3 - 2$, also
$[\Q(\sqrt[3]{2}):\Q] = 3$.
\end{beispiel}

\begin{beispiel}
Der algebraische Abschluss $\overline{\Q}$ von $\Q$ ist algebraisch über $\Q$,
aber von unendlichem Grad. Die reellen Zahlen $\pi$ und $e$ sind transzendent
über $\Q$ (Lindemann-Weierstraß).
\end{beispiel}

\subsection{Einfache Erweiterungen und Kompositum}

\begin{definition}
Das \emph{Kompositum} $L_1 L_2$ zweier Erweiterungen $L_1, L_2$ von $K$ in
einem gemeinsamen Oberkörper ist der kleinste Teilkörper, der $L_1$ und $L_2$
enthält.
\end{definition}

% =============================================================================
\section{Zerfällungskörper und algebraischer Abschluss}
% =============================================================================

\subsection{Zerfällungskörper}

\begin{definition}
Ein \emph{Zerfällungskörper} von $f \in K[x]$ ist eine Körpererweiterung
$L/K$, so dass $f$ über $L$ in Linearfaktoren zerfällt, d.\,h.
$f(x) = a \prod_{i=1}^n (x - \alpha_i)$ mit $\alpha_i \in L$, und
$L = K(\alpha_1,\dots,\alpha_n)$.
\end{definition}

\begin{theorem}[Existenz und Eindeutigkeit von Zerfällungskörpern]
\label{thm:zerfall}
Für jedes $f \in K[x]$ existiert ein Zerfällungskörper $L/K$. Dieser ist bis
auf $K$-Isomorphie eindeutig, und es gilt $[L:K] \leq (\deg f)!$.
\end{theorem}

\begin{proof}[Beweisskizze]
\emph{Existenz}: Durch Induktion über $\deg f$. Hat $f$ keinen irreduziblen
Faktor vom Grad $> 1$ in $K[x]$, ist $K$ selbst der Zerfällungskörper. Andernfalls
wähle einen irreduziblen Faktor $p \mid f$, setze $K_1 = K[x]/(p)$, adjungiere
eine Wurzel $\alpha_1$ von $p$, und fahre mit $f/(x-\alpha_1)$ über $K_1$ fort.

\emph{Eindeutigkeit}: Zwei Zerfällungskörper $L, L'$ sind durch einen
$K$-Isomorphismus verbunden, der schrittweise entlang der Wurzeln konstruiert wird.
\end{proof}

\begin{beispiel}
Der Zerfällungskörper von $x^3 - 2$ über $\Q$ ist
$\Q(\sqrt[3]{2}, \omega)$ mit $\omega = e^{2\pi i/3}$.
Es gilt $[\Q(\sqrt[3]{2},\omega):\Q] = 6$.
\end{beispiel}

\subsection{Algebraischer Abschluss}

\begin{definition}
Ein Körper $\Omega$ heißt \emph{algebraisch abgeschlossen}, wenn jedes
nicht-konstante Polynom in $\Omega[x]$ eine Wurzel in $\Omega$ besitzt.
Der \emph{algebraische Abschluss} $\overline{K}$ von $K$ ist der kleinste
algebraisch abgeschlossene Körper, der $K$ enthält.
\end{definition}

\begin{theorem}[Existenz des algebraischen Abschlusses (Steinitz, 1910)]
Jeder Körper $K$ besitzt einen algebraischen Abschluss $\overline{K}$,
eindeutig bis auf $K$-Isomorphie.
\end{theorem}

% =============================================================================
\section{Separabilität und Normalität}
% =============================================================================

\subsection{Separable Erweiterungen}

\begin{definition}
Ein Polynom $f \in K[x]$ heißt \emph{separabel}, wenn es keine mehrfachen
Wurzeln in $\overline{K}$ besitzt, d.\,h. $\gcd(f, f') = 1$ in $K[x]$.
Ein algebraisches Element $\alpha$ über $K$ heißt separabel, wenn
$\mathrm{min}_K(\alpha)$ separabel ist. Eine algebraische Erweiterung $L/K$
heißt \emph{separabel}, wenn jedes Element von $L$ separabel über $K$ ist.
\end{definition}

\begin{bemerkung}
Jede algebraische Erweiterung eines Körpers der Charakteristik $0$ (z.\,B.\ $\Q$)
oder eines endlichen Körpers ist separabel. Nicht-separable Erweiterungen treten
nur in Charakteristik $p > 0$ auf.
\end{bemerkung}

\begin{theorem}[Satz vom primitiven Element]
Ist $L/K$ eine endliche separable Erweiterung, so gilt $L = K(\theta)$ für
ein $\theta \in L$.
\end{theorem}

\subsection{Normale Erweiterungen}

\begin{definition}
Eine endliche Erweiterung $L/K$ heißt \emph{normal}, wenn sie der
Zerfällungskörper eines Polynoms $f \in K[x]$ ist. Äquivalent: Für jedes
irreduzible $p \in K[x]$ gilt: Hat $p$ eine Wurzel in $L$, so zerfällt $p$
vollständig in $L$.
\end{definition}

\begin{beispiel}
$\Q(\sqrt[3]{2})/\Q$ ist \emph{nicht} normal: $x^3 - 2$ hat eine Wurzel
$\sqrt[3]{2} \in \Q(\sqrt[3]{2})$, zerfällt dort aber nicht vollständig.
Der Zerfällungskörper $\Q(\sqrt[3]{2},\omega)$ ist normal über $\Q$.
\end{beispiel}

% =============================================================================
\section{Die Galois-Gruppe}
\label{sec:galoisgruppe}
% =============================================================================

\begin{definition}
Sei $L/K$ eine Körpererweiterung. Die \emph{Galois-Gruppe} ist
\[
  \Gal(L/K) = \Aut_K(L) = \{\sigma: L \xrightarrow{\sim} L \mid \sigma|_K = \mathrm{id}_K\}.
\]
Dies ist eine Gruppe unter Komposition von Abbildungen.
\end{definition}

\begin{definition}
Eine endliche Erweiterung $L/K$ heißt \emph{Galois-Erweiterung}, wenn sie
sowohl normal als auch separabel ist. In diesem Fall gilt $|\Gal(L/K)| = [L:K]$.
\end{definition}

\begin{beispiel}[Kreisteilungserweiterung]
Sei $\zeta_n = e^{2\pi i/n}$. Die Erweiterung $\Q(\zeta_n)/\Q$ ist galoissch
mit
\[
  \Gal(\Q(\zeta_n)/\Q) \cong (\Z/n\Z)^\times.
\]
Jeder Automorphismus $\sigma_k$ (mit $\ggT(k,n)=1$) sendet $\zeta_n \mapsto \zeta_n^k$.
\end{beispiel}

\begin{beispiel}[Galois-Gruppe von $x^3 - 2$]
Der Zerfällungskörper $L = \Q(\sqrt[3]{2}, \omega)$ über $\Q$ hat $[L:\Q] = 6$.
Die Galois-Gruppe $\Gal(L/\Q) \cong S_3$, die symmetrische Gruppe auf drei
Elementen, die durch Permutation der drei Wurzeln
$\{\sqrt[3]{2}, \omega\sqrt[3]{2}, \omega^2\sqrt[3]{2}\}$ operiert.
\end{beispiel}

\begin{beispiel}[Galois-Gruppe einer quadratischen Gleichung]
Für $f(x) = x^2 - d$ mit $d \in \Q$ kein vollständiges Quadrat gilt
$\Gal(\Q(\sqrt{d})/\Q) \cong \Z/2\Z$,
erzeugt durch $\sqrt{d} \mapsto -\sqrt{d}$.
\end{beispiel}

\subsection{Diskriminante und Galois-Gruppe}

\begin{theorem}
Sei $f \in K[x]$ separabel vom Grad $n$ mit Zerfällungskörper $L/K$.
Die Galois-Gruppe $G = \Gal(L/K)$ bettet sich in $S_n$ ein via ihre Wirkung
auf die Wurzeln. Ist $\Delta = \prod_{i < j}(\alpha_i - \alpha_j)^2$ die
Diskriminante:
\begin{enumerate}[label=(\roman*)]
  \item $G \subseteq A_n$ genau dann, wenn $\sqrt{\Delta} \in K$.
  \item Für $n = 3$: $G \cong A_3 \cong \Z/3\Z$ genau dann, wenn
        $\sqrt{\Delta} \in K$; sonst $G \cong S_3$.
  \item Für $n = 2$: $G \cong \Z/2\Z$ genau dann, wenn
        $\Delta \notin K^{\times 2}$.
\end{enumerate}
\end{theorem}

\subsection{Frobenius-Endomorphismus}

\begin{definition}
Für einen endlichen Körper $\F_{p^n}$ ist der \emph{Frobenius-Automorphismus}
$\varphi: x \mapsto x^p$. Es gilt
\[
  \Gal(\F_{p^n}/\F_p) = \langle \varphi \rangle \cong \Z/n\Z.
\]
\end{definition}

% =============================================================================
\section{Der Hauptsatz der Galois-Theorie}
% =============================================================================

\begin{theorem}[Hauptsatz der Galois-Theorie]
\label{thm:hauptsatz}
Sei $L/K$ eine endliche Galois-Erweiterung mit Galois-Gruppe $G = \Gal(L/K)$.
Es gibt eine ordnungsumkehrende Bijektion
\[
  \left\{\text{Untergruppen } H \leq G\right\}
  \longleftrightarrow
  \left\{\text{Zwischenkörper } K \subseteq F \subseteq L\right\}
\]
gegeben durch $H \mapsto L^H = \Fix(H) = \{x \in L : \sigma(x) = x\ \forall\,\sigma \in H\}$
und $F \mapsto \Gal(L/F)$. Außerdem gilt:
\begin{enumerate}[label=(\roman*)]
  \item $[L:L^H] = |H|$ und $[L^H:K] = [G:H]$.
  \item $H$ ist normal in $G$ genau dann, wenn $L^H/K$ eine Galois-Erweiterung
        ist; in diesem Fall $\Gal(L^H/K) \cong G/H$.
  \item Die triviale Untergruppe $\{e\}$ entspricht $L$, und $G$ selbst
        entspricht $K$.
\end{enumerate}
\end{theorem}

\begin{proof}[Beweisskizze]
Die wesentlichen Hilfsmittel sind:
\begin{itemize}
  \item \emph{Artins Satz}: Ist $G$ eine endliche Automorphismengruppe von $L$
        mit Fixkörper $K = L^G$, so gilt $[L:K] = |G|$, und $L/K$ ist galoissch
        mit Gruppe $G$.
  \item \emph{Injektivität}: Die Abbildung $H \mapsto L^H$ ist injektiv, da
        jedes $\sigma \in H \setminus \{e\}$ ein Element von $L$ bewegt, das
        nicht von allen $\sigma \in H$ festgelassen wird.
  \item \emph{Surjektivität}: Für jeden Zwischenkörper $F$ gilt
        $F = L^{\Gal(L/F)}$ nach Artins Satz, angewandt auf $\Gal(L/F)$
        mit Fixkörper $F$.
  \item Die Gradformel $[L:L^H] = |H|$ folgt direkt aus Artins Satz.
\end{itemize}
\end{proof}

\begin{beispiel}[Galois-Korrespondenz für $\Q(\sqrt{2},\sqrt{3})$]
Sei $L = \Q(\sqrt{2},\sqrt{3})$. Dann $[L:\Q] = 4$ und
$G = \Gal(L/\Q) \cong \Z/2\Z \times \Z/2\Z = \{e, \sigma, \tau, \sigma\tau\}$
mit $\sigma(\sqrt{2}) = -\sqrt{2}$, $\sigma(\sqrt{3}) = \sqrt{3}$,
$\tau(\sqrt{2}) = \sqrt{2}$, $\tau(\sqrt{3}) = -\sqrt{3}$.

Untergruppen und zugehörige Fixkörper:
\begin{center}
\begin{tabular}{ll}
\toprule
Untergruppe $H$ & Fixkörper $L^H$ \\
\midrule
$\{e\}$ & $L = \Q(\sqrt{2},\sqrt{3})$ \\
$\{e, \sigma\}$ & $\Q(\sqrt{3})$ \\
$\{e, \tau\}$ & $\Q(\sqrt{2})$ \\
$\{e, \sigma\tau\}$ & $\Q(\sqrt{6})$ \\
$G$ & $\Q$ \\
\bottomrule
\end{tabular}
\end{center}
Da $G$ abelsch ist, sind alle Untergruppen normal, d.\,h. alle
Zwischenerweiterungen sind galoissch über $\Q$.
\end{beispiel}

% =============================================================================
\section{Lösbarkeit durch Radikale}
% =============================================================================

\subsection{Radikalerweiterungen}

\begin{definition}
Eine Erweiterung $L/K$ heißt \emph{Radikalerweiterung}, wenn ein Turm
\[
  K = K_0 \subset K_1 \subset \cdots \subset K_r = L
\]
existiert mit $K_{i+1} = K_i(\alpha_i)$ und $\alpha_i^{n_i} \in K_i$ für
gewisse $n_i \in \Z_{>0}$. Ein Polynom $f \in K[x]$ ist \emph{durch Radikale
lösbar}, wenn sein Zerfällungskörper in einer Radikalerweiterung von $K$
enthalten ist.
\end{definition}

\begin{definition}
Eine Gruppe $G$ heißt \emph{auflösbar}, wenn sie eine subnormale Reihe
\[
  \{e\} = G_r \trianglelefteq G_{r-1} \trianglelefteq \cdots \trianglelefteq G_0 = G
\]
mit abelschen Faktoren $G_{i-1}/G_i$ besitzt.
\end{definition}

\begin{beispiel}[Auflösbare Gruppen kleiner Ordnung]
$S_1, S_2, S_3, S_4$ sind auflösbar:
\begin{align*}
  S_3 &\supset A_3 \cong \Z/3\Z \supset \{e\},
  \quad \text{Faktoren: } \Z/2\Z,\ \Z/3\Z \\
  S_4 &\supset A_4 \supset V_4 \cong (\Z/2\Z)^2 \supset \Z/2\Z \supset \{e\},
  \quad \text{alle Faktoren abelsch.}
\end{align*}
\end{beispiel}

\subsection{Der Hauptsatz über Lösbarkeit durch Radikale}

\begin{theorem}[Galois, 1832]
\label{thm:loesbar}
Sei $K$ ein Körper der Charakteristik $0$ und $f \in K[x]$. Dann ist $f$
durch Radikale über $K$ lösbar genau dann, wenn $\Gal(f/K)$ eine auflösbare
Gruppe ist.
\end{theorem}

\begin{proof}[Beweisskizze]
($\Rightarrow$) Liegt $L$ in einem Radikalturm $K = K_0 \subset \cdots \subset
K_r$, so kann man jeden Schritt durch eine zyklische Erweiterung ersetzen
(nach Adjunktion geeigneter Einheitswurzeln), was eine Kette von Galois-Gruppen
mit zyklischen — also abelschen — Quotienten liefert. Damit ist $\Gal(L/K)$
auflösbar.

($\Leftarrow$) Ist $G = \Gal(L/K)$ auflösbar mit der Reihe
$\{e\} = G_r \trianglelefteq \cdots \trianglelefteq G_0 = G$, so bilden die
Fixkörper $K_i = L^{G_i}$ Zwischenerweiterungen mit
$\Gal(K_{i+1}/K_i) \cong G_i/G_{i+1}$ zyklisch. Nach der Kummer-Theorie
(nach Adjunktion geeigneter Einheitswurzeln) ist jede zyklische Erweiterung
vom Grad $n$ von der Form $K_i(\alpha_i)$ mit $\alpha_i^n \in K_i$, also
radikal.
\end{proof}

\subsection{Explizite Formeln für Grad 2, 3 und 4}

Der Vollständigkeit halber führen wir die klassischen Radikalformeln auf.

\begin{beispiel}[Grad 2]
$x^2 + px + q = 0$ hat die Lösungen
$x = \dfrac{-p \pm \sqrt{p^2 - 4q}}{2}$.
Die Galois-Gruppe ist trivial, wenn $p^2 - 4q$ ein Quadrat in $K$ ist,
sonst $\Z/2\Z$; in beiden Fällen auflösbar.
\end{beispiel}

\begin{beispiel}[Grad 3: Cardano-Formel]
Für $x^3 + px + q = 0$ (deprimierte Form) sei $\Delta = -4p^3 - 27q^2$.
Die Wurzeln sind
\[
  x_k = \omega^k \sqrt[3]{-\frac{q}{2} + \sqrt{-\frac{\Delta}{108}}}
       + \omega^{-k} \sqrt[3]{-\frac{q}{2} - \sqrt{-\frac{\Delta}{108}}},
  \quad k = 0,1,2,
\]
mit $\omega = e^{2\pi i/3}$. Der Radikalturm ist
$\Q \subset \Q(\omega) \subset \Q(\omega, u)$, wobei $u$ eine Kubikwurzel ist.
\end{beispiel}

\begin{beispiel}[Grad 4: Ferraris Methode]
Ferrari (1545) reduzierte die Quartik durch eine Hilfsgröße (\emph{Resolventenkubik}).
Für $f(x) = x^4 + bx^2 + cx + d$ lautet die kubische Resolvente
$y^3 - by^2 - 4dy + (4bd - c^2)$. Nach Lösung der Kubik durch Cardano
reduziert sich die Quartik auf zwei quadratische Gleichungen, was einen
Radikalturm der Tiefe 4 ergibt.
\end{beispiel}

% =============================================================================
\section{Die Unlösbarkeit der Quintik}
% =============================================================================

\subsection{Einfachheit von $A_5$}

\begin{theorem}[$A_5$ ist einfach]
\label{thm:a5einfach}
Die alternierende Gruppe $A_5$ der geraden Permutationen auf 5 Elementen
ist einfach, d.\,h. sie besitzt keinen echten nicht-trivialen Normalteiler.
\end{theorem}

\begin{proof}
$|A_5| = 60$. Die Konjugationsklassen in $A_5$ haben die Größen
$1, 12, 12, 15, 20$. Ein Normalteiler $N \trianglelefteq A_5$ ist eine
Vereinigung von Konjugationsklassen, die die Identität enthält; also ist $|N|$
ein Teiler von $60$ und muss als Teilsumme aus $\{1, 12, 12, 15, 20\}$
(die $1$ enthaltend) darstellbar sein. Die einzigen solchen Teilsummen, die
$60$ teilen, sind $1$ und $60$, also gilt $N = \{e\}$ oder $N = A_5$.
\end{proof}

\subsection{Nicht-Auflösbarkeit von $S_5$}

\begin{theorem}[$S_5$ ist nicht auflösbar]
\label{thm:s5nichtloesbar}
Die symmetrische Gruppe $S_5$ ist nicht auflösbar.
\end{theorem}

\begin{proof}
Die einzigen Normalteiler von $S_5$ sind $\{e\}$, $A_5$ und $S_5$. Die Reihe
$S_5 \supset A_5 \supset \{e\}$ hat den Faktor $A_5/\{e\} \cong A_5$, der
nicht abelsch ist ($|A_5| = 60$, und eine einfache nicht-triviale Gruppe
der Ordnung $> 1$ ist nicht abelsch, es sei denn zyklisch). Also ist $S_5$
nicht auflösbar.
\end{proof}

\subsection{Der Satz von Abel-Ruffini-Galois}

\begin{theorem}[Abel-Ruffini, 1824; Galois, 1832]
\label{thm:abelruffini}
Es gibt keine allgemeine Lösungsformel durch Radikale für Polynomgleichungen
vom Grad $n \geq 5$ über $\Q$.
\end{theorem}

\begin{proof}
Es genügt, ein Polynom vom Grad $5$ mit Galois-Gruppe $S_5$ anzugeben.
Betrachte $f(x) = x^5 - x - 1 \in \Q[x]$. Man kann zeigen (mit Hilfe der
Diskriminante und Dedekinds Reduktionslemma modulo Primzahlen), dass
$\Gal(f/\Q) \cong S_5$. Nach Satz~\ref{thm:loesbar} ist $f$ durch Radikale
lösbar genau dann, wenn $S_5$ auflösbar ist. Da $S_5$ nicht auflösbar ist
(Satz~\ref{thm:s5nichtloesbar}), besitzt $f$ keine Darstellung seiner Wurzeln
durch Radikale und Koeffizienten.

Da das \emph{generische} Grad-$n$-Polynom über $\Q(t_1,\dots,t_n)$ die
Galois-Gruppe $S_n$ hat und $S_n$ für $n \geq 5$ nicht auflösbar ist,
existiert keine allgemeine Radikalformel für beliebige Polynome vom Grad
$n \geq 5$.
\end{proof}

\begin{bemerkung}
Der Satz besagt \emph{nicht}, dass kein einziges Grad-5-Polynom durch
Radikale lösbar wäre. Viele sind es durchaus: z.\,B.\ $x^5 - 2$ ist lösbar
(seine Galois-Gruppe ist die Frobenius-Gruppe $\Z/5\Z \rtimes \Z/4\Z$ der
Ordnung 20, die auflösbar ist). Die Aussage ist, dass keine
\emph{universelle Formel} in den Koeffizienten für alle Polynome funktioniert.
\end{bemerkung}

\begin{bemerkung}
Kompositionsreihen nach Grad:
\begin{center}
\begin{tabular}{lccc}
\toprule
Grad & Galois-Gruppe & Auflösbar? & Formel \\
\midrule
1 & $\{e\}$ & Ja & $x = -b/a$ \\
2 & $\Z/2\Z$ & Ja & Quadratische Formel \\
3 & $\subseteq S_3$ & Ja & Cardano (1545) \\
4 & $\subseteq S_4$ & Ja & Ferrari (1545) \\
$\geq 5$ & kann $S_n$ sein & Nein & Keine allgemeine Formel \\
\bottomrule
\end{tabular}
\end{center}
\end{bemerkung}

% =============================================================================
\section{Der Kronecker-Weber-Satz}
% =============================================================================

\subsection{Kreisteilungskörper}

\begin{definition}
Das $n$-te \emph{Kreisteilungspolynom} ist
\[
  \Phi_n(x) = \prod_{\substack{1 \leq k \leq n \\ \ggT(k,n) = 1}} (x - e^{2\pi i k/n})
  \in \Z[x],
\]
vom Grad $\varphi(n) = |(\Z/n\Z)^\times|$. Der $n$-te Kreisteilungskörper ist
$\Q(\zeta_n)$ mit $\zeta_n = e^{2\pi i/n}$.
\end{definition}

\begin{theorem}[Galois-Gruppe des Kreisteilungskörpers]
$\Q(\zeta_n)/\Q$ ist eine Galois-Erweiterung vom Grad $\varphi(n)$ mit
$\Gal(\Q(\zeta_n)/\Q) \cong (\Z/n\Z)^\times$
(abelsche Gruppe).
\end{theorem}

\subsection{Der Satz}

\begin{theorem}[Kronecker-Weber, 1853/1886]
\label{thm:kronweber}
Jede endliche abelsche Erweiterung von $\Q$ ist in einem Kreisteilungskörper
enthalten. Genauer: Ist $L/\Q$ eine endliche Galois-Erweiterung mit abelscher
Galois-Gruppe, so gilt $L \subseteq \Q(\zeta_n)$ für ein $n \in \Z_{>0}$.
\end{theorem}

\begin{bemerkung}
Kronecker formulierte den Satz 1853, jedoch ohne vollständigen Beweis.
Weber gab 1886 einen (fast vollständigen) Beweis. Den ersten vollständigen
Beweis lieferte Hilbert (1896).
\end{bemerkung}

\begin{proof}[Beweisskizze via lokale-globale Argumente]
Der moderne Beweis verläuft in zwei Schritten:
\begin{enumerate}
  \item \emph{Lokaler Kronecker-Weber}: Jede endliche abelsche Erweiterung von
        $\Q_p$ (den $p$-adischen rationalen Zahlen) ist in einem Kreisteilungskörper
        über $\Q_p$ enthalten (bewiesen über die Struktur von $\Q_p^\times$ mittels
        lokaler Klassenkörpertheorie oder direkten Argumenten).
  \item \emph{Globale Zusammensetzung}: Eine endliche abelsche Erweiterung
        $L/\Q$ ist durch ihre Vervollständigungen an allen Primzahlen bestimmt.
        Aus dem lokalen Ergebnis an jeder Primzahl $p \mid [L:\Q]$ konstruiert
        man ein $n$ mit $L \subseteq \Q(\zeta_n)$.
\end{enumerate}
\end{proof}

\begin{beispiel}
Die Erweiterung $\Q(\sqrt{d})/\Q$ für quadratfreies $d$ hat Galois-Gruppe
$\Z/2\Z$ (abelsch). Nach Kronecker-Weber gilt
$\Q(\sqrt{d}) \subseteq \Q(\zeta_n)$ für ein geeignetes $n$. Tatsächlich:
$\Q(\sqrt{d}) \subseteq \Q(\zeta_{|d|})$ (bzw.\ $\Q(\zeta_{4|d|})$ für
$d \equiv 2,3 \pmod{4}$).
\end{beispiel}

\begin{beispiel}
$\Q(i) = \Q(\sqrt{-1}) \subseteq \Q(\zeta_4)$, da $\zeta_4 = i$.
\end{beispiel}

\begin{bemerkung}
Kronecker-Weber ist der erste große Satz der Klassenkörpertheorie. Seine
Verallgemeinerung — das Artinsche Reziprozitätsgesetz — beschreibt abelsche
Erweiterungen eines beliebigen Zahlkörpers $K$ mittels der Idelklassengruppe
von $K$.
\end{bemerkung}

% =============================================================================
\section{Konstruierbarkeit mit Zirkel und Lineal}
% =============================================================================

\subsection{Algebraisches Kriterium}

\begin{definition}
Eine reelle Zahl $\alpha$ heißt \emph{konstruierbar}, wenn sie aus $0$ und $1$
durch endlich viele Körperoperationen und Quadratwurzeln gewonnen werden kann.
Äquivalent: $\alpha$ ist konstruierbar genau dann, wenn ein Turm
\[
  \Q = F_0 \subset F_1 \subset \cdots \subset F_m
\]
mit $[F_{i+1}:F_i] = 2$ für alle $i$ und $\alpha \in F_m$ existiert.
\end{definition}

\begin{theorem}[Algebraisches Kriterium für Konstruierbarkeit]
\label{thm:konstruierbar}
Eine reelle Zahl $\alpha$ ist konstruierbar genau dann, wenn $[\Q(\alpha):\Q]$
eine Zweierpotenz ist.
\end{theorem}

\subsection{Klassische Unlösbarkeitsergebnisse}

\begin{korollar}[Würfelverdopplung — unmöglich]
$\sqrt[3]{2}$ ist mit Zirkel und Lineal nicht konstruierbar.
\end{korollar}

\begin{proof}
$[\Q(\sqrt[3]{2}):\Q] = 3$, keine Zweierpotenz.
\end{proof}

\begin{korollar}[Winkeldreiteilung — im Allgemeinen unmöglich]
Ein Winkel von $60°$ kann im Allgemeinen nicht mit Zirkel und Lineal
dreigestrahlt werden.
\end{korollar}

\begin{proof}
Die Dreiteilung von $60°$ erfordert die Konstruktion von $\cos 20°$, das die
Gleichung $8x^3 - 6x - 1 = 0$ erfüllt (irreduzibel über $\Q$, Grad 3). Daher
$[\Q(\cos 20°):\Q] = 3$, keine Zweierpotenz.
\end{proof}

\begin{korollar}[Quadratur des Kreises — unmöglich]
Die reelle Zahl $\sqrt{\pi}$ ist transzendent (Lindemann, 1882), also nicht
algebraisch und insbesondere nicht konstruierbar.
\end{korollar}

\subsection{Satz von Gauß-Wantzel: Reguläre Polygone}

\begin{definition}
Die \emph{Fermat-Primzahlen} sind Primzahlen der Form $F_k = 2^{2^k} + 1$.
Die bekannten Fermat-Primzahlen sind $3, 5, 17, 257, 65537$
(für $k = 0,1,2,3,4$).
\end{definition}

\begin{theorem}[Gauß-Wantzel, 1801/1837]
\label{thm:gausswantzel}
Ein reguläres $n$-Eck ist mit Zirkel und Lineal konstruierbar genau dann, wenn
\[
  n = 2^a \cdot p_1 \cdot p_2 \cdots p_k,
\]
wobei $a \geq 0$ und $p_1 < p_2 < \cdots < p_k$ paarweise verschiedene
Fermat-Primzahlen sind.
\end{theorem}

\begin{proof}[Beweisskizze]
Die Konstruktion eines regulären $n$-Ecks ist äquivalent zur Konstruktion von
$\zeta_n = e^{2\pi i/n}$, das in $\Q(\zeta_n)$ mit $[\Q(\zeta_n):\Q] = \varphi(n)$
liegt. Nach Satz~\ref{thm:konstruierbar} ist dies möglich genau dann, wenn
$\varphi(n)$ eine Zweierpotenz ist. Die elementare Zahlentheorie zeigt, dass
$\varphi(n) = 2^a$ genau dann gilt, wenn $n$ die angegebene Form hat.
\end{proof}

\begin{beispiel}[Heptadekagon]
Gauß entdeckte 1796 (im Alter von 19 Jahren), dass das reguläre 17-Eck
($n = 17 = F_2$) konstruierbar ist, da $\varphi(17) = 16 = 2^4$. Die explizite
Konstruktion beinhaltet $\cos(2\pi/17)$, das ausgedrückt werden kann als
\[
  \cos\!\left(\frac{2\pi}{17}\right)
  = \frac{-1 + \sqrt{17} + \sqrt{34 - 2\sqrt{17}}
    + 2\sqrt{17 + 3\sqrt{17} - \sqrt{34-2\sqrt{17}} - 2\sqrt{34+2\sqrt{17}}}}{16}.
\]
Diese Entdeckung soll Gauß so begeistert haben, dass er sich statt für
Philologie für Mathematik entschied.
\end{beispiel}

\begin{beispiel}
Das reguläre Siebeneck ($n = 7$) ist \emph{nicht} konstruierbar:
$\varphi(7) = 6$ ist keine Zweierpotenz.
\end{beispiel}

% =============================================================================
\section{Zusammenfassung und Ausblick}
% =============================================================================

Die Galois-Theorie verknüpft Algebra, Geometrie und Arithmetik durch ein
einziges vereinheitlichendes Prinzip: \emph{Symmetrie bestimmt Lösbarkeit.}
Die hier entwickelten Hauptresultate bilden einen kohärenten Bogen:

\begin{enumerate}
  \item Körpererweiterungen vom Grad $n$ entsprechen Polynomen vom Grad $n$
        (§\ref{sec:koerper}).
  \item Galois-Erweiterungen besitzen eine Symmetriegruppe (§\ref{sec:galoisgruppe}).
  \item Der Hauptsatz schafft eine perfekte Dualität zwischen Untergruppen
        und Zwischenkörpern (Satz~\ref{thm:hauptsatz}).
  \item Ein Polynom ist durch Radikale lösbar genau dann, wenn seine
        Symmetriegruppe auflösbar ist (Satz~\ref{thm:loesbar}).
  \item $S_5$ ist nicht auflösbar — dies beweist den Satz von Abel-Ruffini
        (Satz~\ref{thm:abelruffini}).
  \item Abelsche Erweiterungen von $\Q$ sind kreisteilungsartig
        (Satz~\ref{thm:kronweber}).
  \item Konstruierbarkeit reduziert sich auf Zweierpotenzen als Grad
        (Sätze~\ref{thm:konstruierbar}, \ref{thm:gausswantzel}).
\end{enumerate}

\noindent\textbf{Weiterführende Themen.}
Moderne Erweiterungen der Galois-Theorie umfassen:
\begin{itemize}
  \item \emph{Inverses Galois-Problem}: Ist jede endliche Gruppe als Galois-Gruppe
        über $\Q$ realisierbar? (Im Allgemeinen offen; für auflösbare Gruppen
        durch Schafarevitsch gelöst.)
  \item \emph{Étale Kohomologie und Grothendieck}: Verallgemeinerung auf
        algebraische Geometrie via fundamentale Gruppe
        $\pi_1^{\text{ét}}(\mathrm{Spec}\,K) = \Gal(\overline{K}/K)$.
  \item \emph{Langlands-Programm}: Tiefe Vermutungen über Galois-Darstellungen
        $\rho: \Gal(\overline{\Q}/\Q) \to GL_n(\C)$ und automorphe Formen.
  \item \emph{$p$-adische und Iwasawa-Theorie}: Galois-Gruppen unendlicher
        Türme von Zahlkörpern.
\end{itemize}

\addtocounter{section}{-1}
\refstepcounter{section}
\renewcommand{\thesection}{\Alph{section}}

% =============================================================================
\begin{thebibliography}{99}
% =============================================================================

\bibitem{galois1832}
\'{E}.~Galois,
\textit{Mémoire sur les conditions de résolubilit\'{e} des équations par
radicaux},
Manuskript (1832); erstmals veröffentlicht in \textit{Journal de
Math\'{e}matiques Pures et Appliqu\'{e}es} \textbf{11} (1846), 381--444.

\bibitem{abel1824}
N.H.~Abel,
\textit{Mémoire sur les équations algébriques où on démontre l'impossibilité
de la résolution de l'équation générale du cinquième degré},
Christiania, 1824.

\bibitem{ruffini1799}
P.~Ruffini,
\textit{Teoria generale delle equazioni},
Bologna, 1799.

\bibitem{lang2002}
S.~Lang,
\textit{Algebra}, 3. Aufl.,
Graduate Texts in Mathematics \textbf{211},
Springer, New York, 2002.

\bibitem{neukirch1999}
J.~Neukirch,
\textit{Algebraische Zahlentheorie},
Grundlehren der mathematischen Wissenschaften \textbf{322},
Springer, Berlin, 1992.

\bibitem{stewart2015}
I.~Stewart,
\textit{Galois Theory}, 4. Aufl.,
CRC Press, Boca Raton, 2015.

\bibitem{dummitfoote2004}
D.S.~Dummit und R.M.~Foote,
\textit{Abstract Algebra}, 3. Aufl.,
Wiley, Hoboken, 2004.

\bibitem{cox2004}
D.A.~Cox,
\textit{Galois Theory},
Wiley-Interscience, Hoboken, 2004.

\bibitem{artin1944}
E.~Artin,
\textit{Galois Theory},
Notre Dame Mathematical Lectures \textbf{2},
University of Notre Dame Press, 1944.

\bibitem{serre1979}
J.-P.~Serre,
\textit{Lokale Körper} [Local Fields],
Graduate Texts in Mathematics \textbf{67},
Springer, New York, 1979.

\end{thebibliography}

\end{document}
